\subsection{\texorpdfstring{同态 \(\bigoplus_{i} \mathrm{Hom}_{\mathrm{A}}(\mathrm{M}, \mathrm{N}_{i}) \longrightarrow \mathrm{Hom}_{\mathrm{A}}(\mathrm{M}, \bigoplus_{i} \mathrm{N}_{i})\)}{同态 \textbackslash bigoplus\_\{i\} \textbackslash mathrm\{Hom\}\_\{\textbackslash mathrm\{A\}\}(\textbackslash mathrm\{M\}, \textbackslash mathrm\{N\}\_\{i\}) \textbackslash longrightarrow \textbackslash mathrm\{Hom\}\_\{\textbackslash mathrm\{A\}\}(\textbackslash mathrm\{M\}, \textbackslash bigoplus\_\{i\} \textbackslash mathrm\{N\}\_\{i\})}}

设 A 是环，M 是一个 A-模，\((\mathrm{N}_{i})_{i\in\mathrm{I}}\) 是一族 A-模。对元素 \((u_{i})\)（属于 \(\bigoplus_{i}\operatorname{Hom}_{\mathrm{A}}(\mathrm{M},\mathrm{N}_{i})\)），我们让它对应元素 \(m\mapsto(u_{i}(m))\)（属于 \(\operatorname{Hom}_{\mathrm{A}}(\mathrm{M},\bigoplus_{i}\mathrm{N}_{i})\)）。这样我们就定义了一个典范同态

\[
\varphi : \bigoplus_ {i} \operatorname{Hom} _ {\mathrm{A}} (\mathrm{M}, \mathrm{N} _ {i}) \longrightarrow \operatorname{Hom} _ {\mathrm{A}} \left(\mathrm{M}, \bigoplus_ {i} \mathrm{N} _ {i}\right).
\]

\(\varphi\) 显然是单射。设 u 是 \(\operatorname{Hom}_{\mathrm{A}}(\mathrm{M},\bigoplus_{i}\mathrm{N}_{i})\) 的一个元素。则 u 属于 \(\varphi\) 的像当且仅当使 \(pr_{i}\circ u\neq0\) 的指标 i 组成的集合是有限的。当模 M 有限生成时，这一条件自动满足。

因此，若模 M 有限生成，则同态 \(\varphi\) 是双射。

